\subsection{Compatible local norms in every specified pure-three tower}

\begin{statement}{A.140}{Proved}
{The exact wild-place norm map after a radicand change}
Let \(F/\Q_3\) be finite with \(\mu_3\not\subset F\),
and let \(z\in F^\times\) have normalized valuation
\(h=v_F(z)\) prime to three.
In one algebraic closure choose
\[
 \theta_0=z,\qquad \theta_{j+1}^3=\theta_j,\qquad
                         F_j=F(\theta_j).
\tag{A.140.1}
\]
Then \(F_j/F\) is totally ramified of degree \(3^j\).
Every field norm \(F_j^\times\to F^\times\) and
every transition norm
\(F_{j+1}^\times\to F_j^\times\) is surjective.
The same assertions hold on the actual unit groups.
For every specified \(b\in F^\times\), there is an
actual compatible local family
\[
 b_0=b,\qquad b_j\in F_j^\times,\qquad
 N_{F_{j+1}/F_j}(b_{j+1})=b_j,\qquad
                         v_{F_j}(b_j)=v_F(b).
\tag{A.140.2}
\]
Unit targets admit families of units.
All these assertions also hold in the pro-three
completed multiplicative groups, with their
actual norm maps.

In particular, take a rational \(a\) with \(v_3(a)=1\),
fix \(n\ge1\), and choose
\[
 \alpha_n^{3^n}=a,\qquad
 F=\Q_3(\alpha_n),\qquad v_3(Q)=0,\quad Q\in\Q^\times,
 \qquad a'=aQ^{e3^n},\quad e\in\Z .
\tag{A.140.3}
\]
The old level-\(n\) field is unchanged, since
\(\Q(\alpha_nQ^e)=\Q(\alpha_n)\).
Choose the new roots by
\[
       \theta_j^{3^j}=\alpha_nQ^e,\qquad
                      \theta_{j+1}^3=\theta_j .
\tag{A.140.4}
\]
The completion at three of the actual new
level-\((n+j)\) field is \(F_j=\Q_3(\theta_j)\).
Every prescribed old three-adic norm target,
including a specified completed unit, has a
compatible local norm family in these precise
new fields.
This supplies the local map at three; it makes
no assertion that the local family is the
localization of a compatible global family.
\end{statement}
\begin{proof}
\emph{The degrees, total ramification, and the literal Galois closure.}
The valuation of \(\theta_j\) relative to \(F\)
is \(h/3^j\), whose denominator is \(3^j\).
Thus its ramification index over \(F\) is at
least \(3^j\), while its defining polynomial
has degree \(3^j\).
Both inequalities are equalities.
This proves irreducibility, total ramification,
and \(v_{F_j}(\theta_j)=h\).
The degree \([F_j:F]\) is odd, so \(F_j\)
does not contain \(\mu_3\).

For a fixed \(j\ge1\), put
\[
 R_j=F(\mu_{3^j}),\qquad
 \Omega_j=R_j(\theta_j),\qquad
 H_j=\operatorname{Gal}(R_j/F).
\tag{A.140.5}
\]
The extension \(R_j/F\) is abelian.
The cubic root field \(F(z^{1/3})/F\) has
degree three and is not Galois: a Galois
cubic root field would contain \(\mu_3\),
which has degree two over \(F\).
It cannot be a subfield of an abelian
extension of \(F\).
Therefore \(z\) is not a cube in \(R_j\).
Its Kummer class has exact order \(3^j\)
modulo \((R_j^\times)^{3^j}\), so
\([\Omega_j:R_j]=3^j\).
For completeness, a smaller order would
write \(z^{3^{j-1}}=w^{3^j}\), hence
\(z=w^3\xi\) with
\(\xi\in\mu_{3^{j-1}}\); that \(\xi\)
is itself a cube in \(R_j\), a contradiction.

The degree equality and the minimal polynomial
over \(R_j\) give the actual automorphisms
\[
 \begin{gathered}
 \tau(\theta_j)=\zeta_{3^j}\theta_j,\qquad
 \tau|_{R_j}=1,\\
 \sigma_u(\theta_j)=\theta_j,\qquad
 \sigma_u(\zeta_{3^j})=\zeta_{3^j}^{\,u},
                  \quad u\in H_j\subset(\Z/3^j)^\times,\\
 G_j:=\operatorname{Gal}(\Omega_j/F)
       =\langle\tau\rangle\rtimes H_j,\qquad
             \sigma_u\tau\sigma_u^{-1}=\tau^u .
 \end{gathered}
\tag{A.140.6}
\]
Here \(\Omega_j\) is the splitting field of
\(X^{3^j}-z\), and the displayed \(H_j\)
is exactly the stabilizer of the chosen
root \(\theta_j\).
The restriction \(H_j\to\operatorname{Gal}
(F(\mu_3)/F)\) is onto.
Thus it contains an element \(u\equiv-1\pmod3\).
Its commutator with \(\tau\) generates the
whole cyclic group \(\langle\tau\rangle\).
Since \(H_j\) is abelian, this proves
\[
 G_j'=\langle\tau\rangle,\qquad
 G_j'H_j=G_j,\qquad
                  \Omega_j^{\,G_j'H_j}=F .
\tag{A.140.7}
\]
The last field is precisely the largest
abelian extension of \(F\) contained in
the pure root field \(F_j=\Omega_j^{H_j}\).

\emph{The actual non-Galois norm comparison.}
We use the norm limitation theorem with
the subgroup in (A.140.7), retaining its
explicit maps
\cite[III, Theorem 3.5]{MilneCFT}.
For the Galois extensions \(\Omega_j/F_j\)
and \(\Omega_j/F\), local Artin reciprocity
gives the commutative norm diagram
\[
 \begin{array}{ccc}
 F_j^\times&\longrightarrow&H_j^{\rm ab}\\
 {\scriptstyle N_{F_j/F}}\downarrow&&
                   \downarrow{\scriptstyle \text{inclusion modulo }G_j'}\\
 F^\times&\longrightarrow&G_j^{\rm ab}.
 \end{array}
\tag{A.140.8}
\]
The top map is onto.
The right map is onto by \(G_j'H_j=G_j\).
Given \(b\in F^\times\), choose
\(c\in F_j^\times\) whose norm has the
same bottom Artin class as \(b\).
The quotient \(b/N_{F_j/F}(c)\) is a norm
from \(\Omega_j\), say \(N_{\Omega_j/F}(d)\).
Norm transitivity gives the literal preimage
\[
       b=N_{F_j/F}\bigl(c\,N_{\Omega_j/F_j}(d)\bigr).
\tag{A.140.9}
\]
This proves surjectivity without replacing
the pure field by its Galois closure.
In particular it uses the actual subgroup
\(G_j'H_j\), not the maximal abelian
subfield of the closure alone.

Apply the same argument with base \(F_j\)
and radicand \(\theta_j\).
This base still has no \(\mu_3\), and its
normalized radicand valuation is the
integer \(h\) prime to three.
It proves surjectivity of every transition
norm in (A.140.1).
For a totally ramified local extension,
normalized valuations satisfy
\[
                 v_{F_j}(N_{F_{j+1}/F_j}w)
                                  =v_{F_{j+1}}(w).
\tag{A.140.10}
\]
Consequently any preimage of a unit is
already a unit.
This proves surjectivity on actual unit
groups as well as on the fields.

\emph{Compatible local fibres, with their topology.}
For a specified \(b\in F^\times\), put
\[
 X_j(b)=\{w\in F_j^\times:
                       N_{F_j/F}(w)=b\}.
\tag{A.140.11}
\]
This set is nonempty by (A.140.9).
If \(m=v_F(b)\), every element of \(X_j(b)\)
has valuation \(m\), by total ramification.
It is a closed subset of the compact annulus
\(\pi_j^m\mathcal O_{F_j}^\times\), where
\(\pi_j\) is any local uniformizer.
For \(w\in X_j(b)\), transition-norm
surjectivity gives \(w'\in F_{j+1}^\times\)
with norm \(w\).
The same valuation formula places it in
the required annulus, and norm transitivity
places it in \(X_{j+1}(b)\).
Thus the maps \(X_{j+1}(b)\to X_j(b)\)
are continuous surjections of nonempty
compact spaces.
Their inverse limit is nonempty, proving
the actual compatible family (A.140.2).
For a unit target all these annuli are
the actual compact unit groups.

The pro-three completion of a finite
extension of \(\Q_3\) is compact: the
valuation factor completes to \(\Z_3\),
and the unit completion is a finitely
generated compact \(\Z_3\)-module with
its possible finite three-primary torsion.
Every finite coefficient quotient is
finite.
Actual norm surjectivity gives surjectivity
on these quotients and then on their
compact inverse limits.
For a completed target, its completed
norm fibres are therefore nonempty compact
spaces, with surjective transition maps
by the same norm-transitivity argument.
Taking their inverse limit proves the
completed assertion.
This reasoning retains the specified
tower and its norm transitions throughout.

\emph{The precise changing-radicand application.}
The polynomial \(X^{3^n}-a\) in (A.140.3)
is Eisenstein over \(\Q_3\).
Its root field has degree \(3^n\), is
totally ramified, and has no \(\mu_3\).
The normalized valuation of \(\alpha_n\)
in \(F\) is one.
Since \(Q\) is a three-adic unit,
\(z=\alpha_nQ^e\) has the same valuation.
The preceding result applies to exactly
the roots in (A.140.4).
Their equations are
\[
          \theta_j^{\,3^{n+j}}=aQ^{e3^n}=a',
       \qquad
          \alpha_n=\theta_j^{3^j}/Q^e .
\tag{A.140.12}
\]
The second identity gives the specified
old-field embedding into every new field.
The new polynomial is still Eisenstein
at three, so its unique completion there
is \(\Q_3(\theta_j)=F_j\).
Thus the norms in (A.140.2) are literally
the localizations of the indicated
new-tower field norms.
All local norm targets are allowed by
this map, while constructing one
compatible global family remains a
separate unresolved step.
\end{proof}

The exact group and valuation computations are retained in
\path{code/wild_pure_tower_norm_certificate.py}
and its output.
The local Artin and compact-fibre arguments above,
rather than a finite numerical test, supply
the infinite-depth surjectivity and compatibility.
